% =============================================================================
% SECTION: Observability in Distributed Systems
% Chapter: Background and Related Work
% =============================================================================
%
% TODO Checklist (verified in codebase/docs):
% [x] Reviewed Q&A.md Q42-Q43 for observability stack details
% [x] Reviewed src/observability/ for metrics, tracing, middleware implementation
% [x] Reviewed src/health/components.py for health check patterns
% [x] Reviewed ops/prometheus/ and ops/grafana/ for monitoring configuration
% [x] Cross-referenced OpenTelemetry and Prometheus documentation
% [x] This section covers conceptual foundations; implementation details in Ch. 11
%
% =============================================================================

\section{Observability in Distributed Systems}
\label{sec:observability-background}

Observability is a critical concern for operating complex distributed systems. This section introduces the conceptual foundations of observability---metrics, logs, and traces---along with the tools and patterns commonly used to implement them. Platform-specific implementation details are presented in Chapter~\ref{ch:observability}.

\subsection{The Three Pillars of Observability}

Observability refers to the ability to understand the internal state of a system by examining its external outputs. In distributed systems, observability is typically achieved through three complementary signals:

\paragraph{Metrics}

Metrics are numerical measurements collected over time, typically aggregated into time series. They answer questions about system behavior at an aggregate level:

\begin{itemize}
    \item How many requests per second is the system handling?
    \item What is the 99th percentile latency?
    \item What percentage of requests are failing?
    \item How much memory/CPU is being used?
\end{itemize}

Metrics are characterized by:
\begin{itemize}
    \item \textbf{Low cardinality}: Metrics are aggregated, so storage grows slowly with traffic.
    \item \textbf{Fixed schema}: Metric names and label sets are predefined.
    \item \textbf{Statistical aggregation}: Counts, gauges, histograms, and summaries.
    \item \textbf{Alerting}: Metrics power alerting rules for anomaly detection.
\end{itemize}

\paragraph{Logs}

Logs are timestamped records of discrete events. They provide detailed information about what happened at specific moments:

\begin{itemize}
    \item Request received with parameters X, Y, Z.
    \item Error occurred: connection timeout to database.
    \item User authenticated successfully.
    \item Background job completed in 3.2 seconds.
\end{itemize}

Logs are characterized by:
\begin{itemize}
    \item \textbf{High cardinality}: Each event can have unique attributes.
    \item \textbf{Flexible schema}: Log content can vary by event type.
    \item \textbf{Full context}: Logs can include detailed error messages, stack traces, and payloads.
    \item \textbf{Forensics}: Logs enable after-the-fact investigation of issues.
\end{itemize}

\paragraph{Traces}

Traces capture the path of a request through a distributed system, showing how different services interact to fulfill a request:

\begin{itemize}
    \item Request entered at API gateway.
    \item API gateway routed to service A.
    \item Service A called service B and database C in parallel.
    \item Response returned after 45ms total.
\end{itemize}

Traces are characterized by:
\begin{itemize}
    \item \textbf{Causality}: Traces show parent-child relationships between operations.
    \item \textbf{Timing}: Each span includes start time and duration.
    \item \textbf{Context propagation}: Trace context is passed between services.
    \item \textbf{Root cause analysis}: Traces help identify which service caused delays or failures.
\end{itemize}

\subsection{Metrics: Prometheus and Beyond}

Prometheus~% TODO: add bibitem for Prometheus
has become the de facto standard for metrics in cloud-native systems, particularly in Kubernetes environments.

\paragraph{Prometheus Architecture}

Prometheus follows a pull-based model:
\begin{itemize}
    \item \textbf{Targets}: Services expose metrics at an HTTP endpoint (typically \texttt{/metrics}).
    \item \textbf{Scraping}: Prometheus periodically fetches metrics from configured targets.
    \item \textbf{Storage}: Metrics are stored in a time-series database.
    \item \textbf{Querying}: PromQL provides a powerful query language for analysis.
    \item \textbf{Alerting}: Alertmanager handles alert routing and notification.
\end{itemize}

\paragraph{Metric Types}

Prometheus defines four metric types:

\begin{itemize}
    \item \textbf{Counter}: A cumulative value that only increases (e.g., total requests).
    \item \textbf{Gauge}: A value that can increase or decrease (e.g., current memory usage).
    \item \textbf{Histogram}: Observations counted into configurable buckets (e.g., request latency distribution).
    \item \textbf{Summary}: Similar to histogram but calculates quantiles client-side.
\end{itemize}

\paragraph{Labels and Cardinality}

Prometheus metrics can have labels (key-value pairs) that add dimensions:

\begin{lstlisting}[language=Python,caption={Example Prometheus metrics with labels}]
# Counter with labels
http_requests_total{method="GET", path="/api/v1/users", status="200"} 12345

# Histogram with labels
http_request_duration_seconds_bucket{method="POST", le="0.5"} 1000
http_request_duration_seconds_bucket{method="POST", le="1.0"} 1500
\end{lstlisting}

High-cardinality labels (e.g., user IDs, request IDs) should be avoided as they create many time series, increasing storage and query costs.

\paragraph{Grafana Visualization}

Grafana~% TODO: add bibitem for Grafana
is commonly paired with Prometheus for visualization:
\begin{itemize}
    \item Dashboard creation with multiple panels.
    \item PromQL integration for querying metrics.
    \item Alerting integration with various notification channels.
    \item Template variables for dynamic dashboards.
\end{itemize}

\paragraph{Metrics in AI Systems}

AI agent platforms require metrics specific to their domain:
\begin{itemize}
    \item LLM invocation counts and latencies.
    \item Token usage (input and output tokens).
    \item Tool invocation success/failure rates.
    \item Agent run completion times and step counts.
    \item Model provider health status.
    \item Rate limit utilization.
\end{itemize}

\subsection{Logging: Structured and Correlated}

Modern logging practices have evolved from unstructured text to structured, machine-parseable formats.

\paragraph{Structured Logging}

Structured logs encode events as key-value pairs or JSON objects:

\begin{lstlisting}[language=Python,caption={Structured log example}]
{
    "timestamp": "2025-01-15T10:30:45.123Z",
    "level": "INFO",
    "message": "Request completed",
    "request_id": "abc-123",
    "method": "POST",
    "path": "/api/v1/agent-runs",
    "status": 201,
    "duration_ms": 156,
    "user_id": "user-456"
}
\end{lstlisting}

Structured logs enable:
\begin{itemize}
    \item Efficient parsing and indexing.
    \item Consistent field names across services.
    \item Faceted search and filtering.
    \item Correlation by request ID or trace ID.
\end{itemize}

\paragraph{Log Levels}

Standard log levels convey severity:
\begin{itemize}
    \item \textbf{DEBUG}: Detailed information for debugging (typically disabled in production).
    \item \textbf{INFO}: Normal operational events.
    \item \textbf{WARNING}: Unexpected but recoverable situations.
    \item \textbf{ERROR}: Failures that affect individual requests.
    \item \textbf{CRITICAL}: Severe failures affecting system operation.
\end{itemize}

\paragraph{Log Correlation}

In distributed systems, correlating logs across services is essential:
\begin{itemize}
    \item \textbf{Request ID}: A unique identifier assigned at the edge and propagated through all services.
    \item \textbf{Trace ID}: A distributed tracing identifier that links logs to traces.
    \item \textbf{Session ID}: Links logs across multiple requests in a user session.
\end{itemize}

\paragraph{Log Management}

Production systems typically use log aggregation solutions:
\begin{itemize}
    \item \textbf{ELK Stack}: Elasticsearch, Logstash, Kibana.
    \item \textbf{Loki}: Prometheus-style logging by Grafana Labs.
    \item \textbf{Cloud services}: AWS CloudWatch, Google Cloud Logging, Azure Monitor.
\end{itemize}

\paragraph{Logging in AI Systems}

AI agent platforms have specific logging considerations:
\begin{itemize}
    \item \textbf{PII handling}: Prompts and responses may contain personal information.
    \item \textbf{Size management}: LLM interactions can be verbose; selective logging is needed.
    \item \textbf{Audit requirements}: Security-relevant events must be logged immutably.
    \item \textbf{Debug visibility}: Reasoning traces help understand agent behavior.
\end{itemize}

\subsection{Distributed Tracing: OpenTelemetry}

OpenTelemetry~% TODO: add bibitem for OpenTelemetry
has emerged as the vendor-neutral standard for distributed tracing (and increasingly for metrics and logs as well).

\paragraph{OpenTelemetry Concepts}

\begin{itemize}
    \item \textbf{Trace}: The complete path of a request through the system.
    \item \textbf{Span}: A single operation within a trace, with start time, duration, and attributes.
    \item \textbf{Context}: Metadata propagated between services (trace ID, span ID, flags).
    \item \textbf{Exporter}: Component that sends telemetry to backends (Jaeger, Zipkin, cloud services).
    \item \textbf{Instrumentation}: Code that creates spans (auto-instrumentation or manual).
\end{itemize}

\paragraph{Context Propagation}

Trace context is propagated through:
\begin{itemize}
    \item HTTP headers (W3C Trace Context standard).
    \item Message queue metadata.
    \item gRPC metadata.
    \item Custom protocols.
\end{itemize}

\paragraph{Sampling}

Not all traces need to be collected:
\begin{itemize}
    \item \textbf{Head-based sampling}: Decision made at trace start (e.g., sample 10\% of traces).
    \item \textbf{Tail-based sampling}: Decision made after trace completes (e.g., keep all error traces).
    \item \textbf{Adaptive sampling}: Sampling rate adjusts based on traffic.
\end{itemize}

Production systems typically sample traces to control costs while ensuring error traces are always captured.

\paragraph{Tracing Backends}

Popular tracing backends include:
\begin{itemize}
    \item \textbf{Jaeger}: Open-source, CNCF-graduated project.
    \item \textbf{Zipkin}: Open-source, originated at Twitter.
    \item \textbf{Tempo}: Grafana Labs' scalable tracing backend.
    \item \textbf{Cloud services}: AWS X-Ray, Google Cloud Trace, Azure Application Insights.
\end{itemize}

\paragraph{Tracing in AI Systems}

AI agent platforms benefit from tracing for:
\begin{itemize}
    \item Visualizing the flow through agent steps.
    \item Identifying slow LLM calls or tool invocations.
    \item Understanding retry and fallback behavior.
    \item Correlating frontend requests to backend operations.
\end{itemize}

\subsection{Health Checks and Probes}

Health checks enable orchestration systems to manage application lifecycle and traffic routing.

\paragraph{Kubernetes Probe Types}

Kubernetes defines three probe types:

\begin{itemize}
    \item \textbf{Liveness probe}: Determines if the container is running. Failure triggers container restart.
    
    \item \textbf{Readiness probe}: Determines if the container is ready to receive traffic. Failure removes the pod from load balancer endpoints.
    
    \item \textbf{Startup probe}: Determines if the application has started. Used for slow-starting applications to prevent premature liveness failures.
\end{itemize}

\paragraph{Probe Implementation}

Probes are typically implemented as:
\begin{itemize}
    \item \textbf{HTTP GET}: Check returns 2xx status code.
    \item \textbf{TCP socket}: Check opens a connection.
    \item \textbf{Command execution}: Check runs a command in the container.
\end{itemize}

\paragraph{Health Check Design}

Effective health checks follow several principles:

\begin{itemize}
    \item \textbf{Liveness should be simple}: Check that the process is responsive, not full dependency health.
    
    \item \textbf{Readiness should check dependencies}: Check that the application can actually handle requests (database connected, caches warm, etc.).
    
    \item \textbf{Avoid cascading failures}: A dependency being down should affect readiness, not liveness.
    
    \item \textbf{Fast response}: Health checks should respond quickly to avoid timeout issues.
    
    \item \textbf{Appropriate failure thresholds}: Configure failure thresholds to avoid flapping.
\end{itemize}

\paragraph{Component Health}

Complex applications may have multiple components with independent health:
\begin{itemize}
    \item Database connectivity.
    \item Cache connectivity.
    \item External API availability.
    \item Background worker status.
    \item Model provider health.
\end{itemize}

Aggregating component health into overall readiness requires policy decisions about which components are critical.

\subsection{Observability Patterns for AI Platforms}

AI agent platforms present specific observability challenges and patterns:

\paragraph{LLM Observability}

Monitoring LLM interactions requires:
\begin{itemize}
    \item \textbf{Latency tracking}: Per-model, per-provider latency distributions.
    \item \textbf{Token counting}: Input and output token usage for cost tracking.
    \item \textbf{Error classification}: Distinguishing rate limits, timeouts, and content errors.
    \item \textbf{Quality metrics}: Optional evaluation of response quality.
\end{itemize}

\paragraph{Agent Observability}

Monitoring agent behavior requires:
\begin{itemize}
    \item \textbf{Run metrics}: Completion rate, step count, duration.
    \item \textbf{Tool metrics}: Invocation counts, success rates, latencies.
    \item \textbf{Reasoning visibility}: Access to agent decision traces.
    \item \textbf{Cost allocation}: Per-run, per-user, per-tenant cost tracking.
\end{itemize}

\paragraph{Multi-Tenant Observability}

Multi-tenant systems need tenant-aware observability:
\begin{itemize}
    \item Metrics labeled by tenant for per-tenant dashboards.
    \item Log filtering by tenant for support investigations.
    \item Trace filtering by tenant for issue diagnosis.
    \item Aggregated and tenant-specific SLO monitoring.
\end{itemize}

\paragraph{Production Readiness Checklist}

A production-ready AI platform should have:

\begin{enumerate}
    \item \textbf{Comprehensive metrics}: Request rates, latencies, error rates, resource usage, business metrics.
    
    \item \textbf{Structured logging}: JSON logs with request correlation, appropriate log levels, PII protection.
    
    \item \textbf{Distributed tracing}: OpenTelemetry instrumentation, sampling configuration, trace-log correlation.
    
    \item \textbf{Health probes}: Kubernetes-compatible liveness, readiness, and startup probes.
    
    \item \textbf{Dashboards}: Grafana dashboards for operational visibility.
    
    \item \textbf{Alerting}: Alert rules for critical conditions with appropriate thresholds and routing.
    
    \item \textbf{Runbooks}: Documentation for responding to common alerts and issues.
\end{enumerate}

The Cineca Agentic Platform implements these observability capabilities, as detailed in Chapter~\ref{ch:observability}, providing operators with comprehensive visibility into system behavior.
